\section{Introduction}

Genetic microarrays, medical imaging, and robotics routinely produce high-dimension, low-sample-size (HDLSS) data, in which the dimension far exceeds the sample size. \citet{hall2005} showed that such data acquire a rigid geometry as the dimension diverges, and \citet{ahn2007} proved that this representation holds under mild conditions; see also \citet{jung2009} and the survey of \citet{aoshima2018}. Two lines of work build on this geometry for classification. One develops linear discriminant rules for many variables \citep{bickel2004,fan2008} and distance-based or geometric classifiers \citep{hall2008,chan2009,aoshima2011,aoshima2014,aoshima2015,aoshima2019}. The other studies margin-based methods, including distance-weighted discrimination \citep{marron2007,qiao2010} and the maximal data-piling direction \citep{ahn2010}.

Within the second line, \citet{nakayama2017} derived the asymptotic behavior of the hard-margin linear support vector machine (SVM) \citep{cortes1995,vapnik2000}. They showed that the SVM is consistent under suitable conditions but can be heavily biased when the classes differ in covariance trace or sample size, and proposed a bias-corrected SVM. \citet{nakayama2020} extended that analysis to the Gaussian kernel.

These results assume that every training point comes from one of the two populations. Real training data can contain mislabeled samples or measurements taken under abnormal conditions \citep{frenay2014}. Robust statistics has a long history \citep{huber1964,hampel1986,rousseeuw1987}, and SVM robustness has been studied through influence functions \citep{christmann2004,steinwart2008}, robust losses \citep{wu2007}, robust optimization \citep{xu2009}, and poisoning attacks \citep{biggio2012}. Those analyses either fix the dimension or use conditions that do not reflect HDLSS geometry.

Suppose one class contains a point whose mean $\boldsymbol\mu_0$ lies far from its nominal mean $\boldsymbol\mu_2$. Can the SVM learn the contrast $\boldsymbol\mu_2-\boldsymbol\mu_0$ instead of the population contrast $\boldsymbol\mu_1-\boldsymbol\mu_2$? Because the hard-margin SVM depends only on support vectors, a single outlier may change the classifier completely. A second question is what happens when both classes are contaminated.

Our contributions are as follows.
\begin{itemize}
  \item We derive the limit of the centered Gram matrix in the presence of outliers (Lemma~\ref{lem:geometry}).
  \item For one-sided contamination, we solve the limiting dual problem in closed form and show that the outliers enter only through one effective distance (Lemma~\ref{lem:reduction} and Theorem~\ref{thm:one-sided}).
  \item We identify a sharp threshold below which the contaminated class loses all inlier support and obtain the corresponding critical number of outliers (Corollary~\ref{cor:allocation}).
  \item We prove convergence of the actual dual value, dual allocation, intercept, and discriminant function, and derive limiting misclassification rates (Theorems~\ref{thm:transfer-value}--\ref{thm:discriminant} and Corollary~\ref{cor:misclass-one}).
  \item For two-sided contamination, we solve the KKT system, identify its multiplier with the limiting intercept, and show that at most one class can fail. Symmetric contamination cancels the intercept shift (Theorem~\ref{thm:two-sided} and Corollaries~\ref{cor:breakdown-two}--\ref{cor:symmetric}).
  \item Seven simulation studies check the finite-dimensional agreement with these limits (Section~\ref{sec:simulations}).
\end{itemize}

All algebraic steps use only linear algebra, probability, and convex optimization \citep{boyd2004}.
